\section{New Thought and Tradition}

\begin{quotex}
Every day, in every way, I am getting better and better
\flright{\textsc{Emile Coué}}
\end{quotex}


Before providing translations of \textbf{Julius Evola}'s essays on the New Thought movement in Europe and the ideas of Emile Coué, I want to offer some personal observations. I will also supplement Evola's material with the comparable ideas from in Fr. \textbf{Alois Wiesinger}'s book \emph{Occult Phenomena}.

Some time ago, at the beginning of the eighties, I received in the mail an offer to subscribe to a glossy magazine promising in depth articles on various traditions. I decided to take a subscription to it. I don't recall the name of it now. They only published a few numbers before they ceased operations. However, they did offer a selection of similar magazines that agreed to fulfill the remainder of the subscription. There was one that grabbed my attention. It represented a point of view that I had never previously confronted. To top it off, it purported to combine two topics of intense interest to me: science and the mind.

The small magazine was called the ``\textbf{Science of Mind}'' and, since it was already paid for, I decided to take it. As I learned more about that movement, I discovered that it was part of a larger movement called ``New Thought'' and was a cousin of Christian Science. Its influence in popular culture has been great as it has appeared in many guises: most recently as the ``Law of Attraction'', but also in Positive Thinking, and even among neo-Christians such as the infamous Joel Osteen. I won't go into the entire history of the movement, but suffice it to say that I spent several years taking various classes and studying the relevant texts.

I always heard in Science of Mind classes that they teach what the ``mystics have always taught''. That motivated me to actually read those mystics for myself. The results were mixed. Yes, there were ideas from the mystics and also many distorted Hermetic and esoteric teachings in New Thought. That just gave me the motivation to seek out the real thing itself. I will here mention a few of those ideas and conclude with what the mystics really taught.

\paragraph{Pragmatism}


The American philosopher \textbf{William James}, in \emph{The Varieties of Religious Experience}, considered New Thought (or Mind Cure) as the religion of healthy mindedness. In contrast, he regarded the major representatives of the major traditions as ``morbid-minded''. This is reminiscent of the flaccid distinction made by some between ``life affirming'' and ``life denying'' religions. Pragmatism does not bother with the truth value of a thought or theory but rather is concerned only with its effects. Specifically, Prof James feels that being ``healthy minded'' in the sense he means by it is a desirable goal. Hence, any belief that leads to that psychological state is considered ``true''.

Pragmatism has infected Western thought in unsuspected ways. It explains, for example, the worst of political correctness. Its adherents, for example, desire a particular social state of affairs. Hence, any idea or outlook that tends to promote that social state is considered true or valid, while any idea that threatens to undermine that state must be vehemently repressed. The result is that anyone who wants to look at things ``scientifically'' or ``objectively'' can only be frustrated in verbal confrontations with the other group. Of course, ``science'' itself can be distorted and used in a pragmatic way, hence it is not the defeater in an argument that some would hope for. Pragmatism is a sophisticated philosophy, so it is beside the point to consider progressives as evil or stupid. A more sophisticated worldview is the antidote.

We see pragmatism used in an unsophisticated way among the so-called new/alt right or neo-traditionalist crowd. The distinguishing feature uniting them is the emphasis on ``race'' or other form of zoological identitarianism. Hence, they evaluate every worldview, politico-philosophical system, religion, or tradition by how well it promotes their view of identity. There is no other standard, not even the truth of the idea. Just to be clear once again, for Tradition, the fundamental identity is that of Spirit, prior to any other form of identity. Such a position is immune to the ravages of pragmatism.

\paragraph{Visualizations}


The goal of New Thought, in all its guises, is some form of magic and techniques associated with magic are thus employed. However, the magic is on the level of the personal. Sorcery, or the use of lower forces, is not deliberately used. Neither is the higher form of sacred magic at issue. New Thought does not believe in a cosmic law or a divine person. Rather divinity is regarded as impersonal, subject to certain laws which an agent can use to his advantage. Oftentimes, the ``universe'' is substituted for divinity, as the ``universe will provide'' or ``respond''.

Hence New Thought movements use some variation of thought control and/or visualization, which is really a form of concentration although I never encountered any teachers teaching concentration as such. The premise is that such practices can alter persons, places, things, or conditions. Visualizing objects of desire will somehow bring them into one's possession. Of course, in Hermetic practice, visualizations are used to lead to higher states of consciousness. For example, one may visualize being in the presence of higher beings, a practice which may lead one to that state. It would be incorrect to believe that such visualizations are ineffective. In actuality, we are creating our selves, or our self images, on an ongoing basis via unconscious visualizations throughout the day.

\paragraph{Mind Treatments and the Unconscious}


Instead of visualizations, usually thoughts are more commonly used. For example, I may say, if I feel ill, that ``I am well and in perfect health''. This is supposed to counteract the initial thought of disease. However, a few minutes per day of positive thoughts are seldom sufficient to counteract the 60,000 thoughts per day that flow randomly through one's consciousness.

A favorite Biblical quote in this movement is ``As a man thinketh in his heart, so he is.'' Now that makes more sense, since the heart is one's intellectual center which, in a fully developed person, dominates the other functions of the mind. But, those thoughts are beyond words and conscious thought are merely their reflection in the mind. As such, they cannot be controlled at will, but through mental purification, perhaps truly higher thoughts will dwell in one's heart. However, such thoughts will be unrelated to the desire for a new Maserati or an Alaskan cruise.

\paragraph{Mantras}


Now there are certainly stories about miraculous events effectuated through mystics. In genuine mysticism, this will be the result of sacred magic and will have a totally different character from personal magic or sorcery. New Thought is not interested in sacred magic, gnosis, or higher states. On the contrary, it prides itself on being ``practical'' (or ``pragmatic''), so that it can use the occult forces for personal benefit in a more or less mechanical way. There is no emphasis at all on any sort of purification of the types explained in Julius Evola's essay recently made available on this blog.

Now the repetition of positive thoughts in a ``mind treatment'' is ultimately ineffective without such purification. Such a mind treatment is really a form of mantra ``although one can repeat it a million times, but as long as it is not \emph{known}, it remains a mere flapping of the lips''. Hence, the missing link in New Thought is that gnosis. Certainly, thought has an influence in the world of manifestation, since knowing the true name of a thing is to know its essence. However, the true gnosis of a mantra requires intense development, which is not a part of any New Thought teaching. Moreover, such a thinking is not merely passive in the sense that the ``universe will respond'' to one's thoughts. Rather, there has to be a force or will associated with it.

\paragraph{The Placebo Effect and the Unconscious}


A simpler explanation of certain results that make New Thought plausible is the ``placebo effect'', by which a materially ineffective medication actual produces a real healing provided the person truly believes the medication is effective. Emile Coué, in his work as a pharmacist, noted that his customers would often report a cure even when he gave them a placebo, or inert sugar pill. Due to space considerations, I will treat this in a separate post.

It is interesting to see how esoteric ideas get distorted in the popular mind. That is why such teachings can be hidden in ``plain view''. It is also why ``initiation'' is necessary, or else their true significance might never be realized.

After a post on the metaphysical basis of what is valid in New Thought, based on Fr. Wiesinger's book \emph{Occult Phenomena}, we will present an alternative worldview, the one that the mystics actually taught.


\flrightit{Posted on 2013-11-18 by Cologero}

\begin{center}* * *\end{center}

\begin{footnotesize}
\begin{sffamily}
\texttt{JA on 2013-11-19 at 12:12 said:}

I am not at all familiar with new thought; is Thelema connected with it ? also this sounds like A Course in Miracles and other things I've found in New Age bookstores.

\hfill

\texttt{Cologero on 2013-11-19 at 13:41 said:}

Yes, JA, it is the pervasive ideology behind New Age movements. A Course in Miracles is certainly related, although Charles Upton gives it a rather extended analysis in one of his books.

\hfill

\texttt{Jacob on 2013-11-21 at 20:02 said:}

I agree Cologero. When I think New Thought I think of people like Napoleon Hill or even Deepak. It seems the New Thought movement seeks shallow and fading things (usually money it seems). It appears to try the exact opposite of Stoicism. Really interested in Evola's thoughts.

\hfill

\texttt{Pickman on 2013-11-27 at 22:40 said:}

Would Echkart Tolle fall into this category? He does seem to abide by the new age theory of human ``spiritual evolution'' yet his practice entails a no-mind state of being similar to the Upanishads.

\hfill

\texttt{Ash on 2013-11-30 at 18:52 said:}

It appears to me that as a whole, New Thought has a lot less publicity now than it did from around the 70's to the early 2000's. A lot of that seems to have been the vehement attack on New Thought by the New Atheist/Skeptic crowd, which did do a lot of good in exposing charlatans who had been fleecing their flocks for quite a pretty penny. I'd say the types of ``thought control'' more present today (at least in cities and academia) are those coming from the ideologies seeking to dismantle ``gender binaries'' and other remaining traditional conceptualizations. Speaking personally, the main appearances New Thought has made around me has been with bored housewives and alienated businessmen.

Of course, that doesn't discount how similar ideas/delusions regarding our minds and how our personal wishes create reality have made inroads into our culture. I highly recommend the documentary ``Century of the Self'', which details this process from Freud til the end of the 20th century, including how these mindsets have been used to exercise control over consumers, political movements, etc. Applying the mental disciplines this blog talks about would disrupt powerful networks in our society if even a tiny percentage of the population began to practice them.

\hfill

\texttt{B on 2014-01-26 at 08:38 said:}

New Thought has shed its skin but you can still see in a more materialist guise neurolinguistic programming (nlp) and many other quirky forms of self-help. If you add a fair dose of pop-culture it reappears as chaos magick, which is often an important component of many magic-centered communities.

Whether `New Thought' is to be blamed for the cultural relativism at the heart of political correctness I am not sure: you could just as well blame Nietzsche or any solipsist for that matter?

For those reading French, here is a link to an article by Guenon on Bergson, and, en passant, William James -- I would imagine this will be relevant:

\url{http://esprit-universel.over-blog.com/ren\%C3\%A9-gu\%C3\%A9non-la-\%C2\%AB-religion-\%C2\%BB-d\%E2\%80\%99un-philosophe}


\end{sffamily}
\end{footnotesize}

